%!TEX root = ../template.tex
%%%%%%%%%%%%%%%%%%%%%%%%%%%%%%%%%%%%%%%%%%%%%%%%%%%%%%%%%%%%%%%%%%%%
%% glossary.tex
%% NOVA thesis document file
%%
%% Glossary definition
%%%%%%%%%%%%%%%%%%%%%%%%%%%%%%%%%%%%%%%%%%%%%%%%%%%%%%%%%%%%%%%%%%%%

\typeout{NT FILE glossary.tex}%

\newglossaryentry{computer}{
  name={computer},
  sort={computer},
  description={An electronic device which is capable of receiving information (data) in a particular form and of performing a sequence of operations in accordance with a predetermined but variable set of procedural instructions (program) to produce a result in the form of information or signals. This is a test that adds a citation~\cite{Artho04} to the glossary!}
}

\newglossaryentry{sdr-control-system}{
  name={\gls{SDR} Control/Interface System},
  sort={sdr-control-system},
  description={The \gls{SDR} Control/Interface System is the main software component developed in this internship. It is responsible for managing the communication between the \gls{SDR} hardware and the software applications that utilize it. It provides an interface for configuring the \gls{SDR} transmission parameters, such as Frequency, \gls{RF}-control, and Sampling Rate, as well as handling data transmission.}
}

\newglossaryentry{rs232-interface}{
  name={\gls{RS232} Interface},
  sort={rs232-interface},
  description={The \gls{RS232} interface is the secondary software component developed in this internship. It is responsible for data encoding and decoding as well as its transmission and reception through the \gls{RS232} physical layer. It is used for intercommunication of pre-existing components.}
}